\scanpage{84}\mspage{71}
% archive sheet number: 83
% Paragraph continues from the preceding sheet.
\uline{toutes} flèches de $\mathrm{Cat}$, et mode de calcul des
fibres \ldots).

4. \uline{Cat. à cofibrations spéciales}

4.1. Je suppose à présent que \textcircled{a} $M$ est stable par
$\underrightarrow{\lim}$ filtrantes, et \textcircled{b} que ses objets sont
accessibles, de sorte qu'on est sous les
conditions d'application du th. de factorisation
(\unsure{X} \S\deleted{\unsure{\S}}1).
% slanted
\margin{\textcircled{c} donnés \inserted{$W$,} $C\subset\mathrm{Fl}\,M$ satisfaisant les conditions Cof 1, Cof 2 de \unsure{p.} 39. Pour Cof 3, cf plus bas.}
\deleted{\unsure{Pos}} On suppose \textcircled{d} $C,TC$\inserted{$=C\cap W$} \uline{stables \unsure{par|pour}
composition transfinie} (c'est le cas si la
structure provient d'une structure de
Quillen close $(W,C,F)$).
\deleted{On suppose de
plus que \textcircled{d}
\[
TF\overset{\mathrm{déf}}{=}C_*\subset W\tag*{(4.1)}
\]
donc $TF\subset W_{\mathrm{univ}}$ (puisqu'il est stable par
chgt. de base). \deleted{\unsure{ce qui est|de surcroît} \unsure{stable|stables} \unsure{on suppose} \unsure{que} \illegible{1}}}
\deleted{\unsure{si on désire, donc}} Posons
\[
\deleted{\unsure{=}}\quad TF\overset{\mathrm{déf}}{=}C_*\subset\mathrm{Fl}\,M\,,
\tag*{(4.1)}
\]
il résulte du th. de factorisation que \deleted{\unsure{toute}}\unsure{tte}
flèche $f$ de $M$ se factorise en
\[
f=pi,\qquad p\in TF,\ i\in C
\tag*{(4.1)}
\]
\uline{pourvu} qu'on fasse l'hypothèse \textcircled{\unsure{e|d}}
\[
\begin{aligned}
&\exists\, C_0\subset C,\quad \text{\uline{petite} partie de }C\text{, telle}\\
&\text{que}\quad C\subset\widetilde{C_0}\quad \text{i.e.}\quad TF=(C_0)_*\,.
\end{aligned}
\tag*{(4.3)}
\]
Sous ces conditions, si on suppose \textcircled{\unsure{f|g}}
\[
TF\subset W\,,
\tag*{(4.4.)}
\]
